% !TeX root = ../../Book2.tex
% 纲要标题：## 附录 A　精确矩阵计算与标准形证书
% 纲要位置：计划纲要.md，第 1742 行。

\chapter{精确矩阵计算与标准形证书}
\label{b2:app:A}

\begin{chapterabstract}
标准形把线性代数的结构问题化为矩阵计算，但一个可靠的输出还应说明工作环、允许的变换以及回到原坐标的方法。本附录补充整数与多项式矩阵的 Smith 算法，给出线性方程与核的坐标公式，讨论一般环的障碍和模素数检验的适用条件，并以有限域、有理数上的精确算例区分矩阵等价与算子相似。
\end{chapterabstract}

\begin{learninggoals}
\begin{itemize}
\item 记录行列变换，核验 Smith 标准形，并用换基矩阵恢复原方程的解与核。
\item 区分结构存在性与算法可执行性，说明复杂度分析还须计入输入长度和系数增长，并辨认非主理想及零因子造成的障碍。
\item 使用模素数秩界和有界同余给出精确证据，辨别近似数据不能确定的代数性质。
\item 对多项式矩阵和有理矩阵写出可以逐项相乘检验的标准形证书。
\end{itemize}
\end{learninggoals}

一个整数矩阵的 Smith 对角形可以很短，但若只报对角元，仍无法知道原方程的未知量怎样变换。在\cref{b2:prop:presentation-matrix}的展示约定中，\(A:R^n\to R^m\) 的目标基对应模的生成元，来源基对应关系生成元。若 \(P\in\operatorname{GL}_m(R)\)、\(Q\in\operatorname{GL}_n(R)\) 且 \(PAQ=D\)，则原目标向量 \(u\) 的标准余核坐标是 \(Pu\)，标准生成元在原坐标中的代表是 \(P^{-1}e_i\)；恢复这两项信息用 \(P\) 即可。\(Q\) 则记录来源坐标的变化：解方程时令 \(x=Qy\)，核中的向量也由 \(Q\) 换回原坐标。因此，同时保存两组矩阵可以核对整个等式，并分别恢复生成元、关系、解与核。

\ProseParagraphBreak

核验一份 Smith 输出时，先检查矩阵等式 \(PAQ=D\)，再检查 \(P,Q\) 在指定工作环上可逆，最后检查 \(D\) 是矩形对角矩阵、非零项先于零项、非零项形成整除链并满足选定的归一化。第一项只验证乘积；第二项才保证余核同构；第三项才使输出成为规定的标准形。例如，整数环上可逆要求行列式为 \(\pm1\)，而域上的一元多项式环中，行列式只须为非零常数。

\ProseParagraphBreak

整数与多项式的带余除法允许逐步压低主元；每次交换行列、加另一行或列的倍数，都要同步记录允许的可逆变换。若改用模素数检验，还须说明它验证的是秩的下界、某个整数矩阵等式，还是换基的可逆性；这些不同结论所需的证据并不相同。

\ProseParagraphBreak

本附录使用\chapref{b2:ch:04}的 Smith 标准形及有理标准形，也需要整数和一元多项式的带余除法、行列式与有限维向量空间的基本计算。存在性与唯一性仍以\cref{b2:thm:smith-normal-form}和\cref{b2:thm:smith-uniqueness}为主要证明位置；这里着重说明如何执行这些论证，以及怎样检查一份具体输出。有限域算例只使用 \(\mathbb F_5\) 的算术，不需要额外的扩域算法。

\ProseParagraphBreak
进一步学习精确算法时，可把本附录的矩阵计算与 Shoup 的《A Computational Introduction to Number Theory and Algebra》\cite[Sections 4.4, 17.3]{Shoup2008NumberTheoryAlgebra}对照：前者介绍利用模运算和整数大小界恢复矩阵乘积，后者给出多项式 Euclid 算法及 Bézout 系数。


\ProseParagraphBreak
\secref{b2:sec:0403}用不变因子分类模，所得抽象同构类型不要求保存换基过程；这里保留换基矩阵，是为了在得到分类以后继续回答原坐标中的问题。

% !TeX root = ../../Book2.tex
% 纲要标题：### A.1 整数与一元多项式矩阵的 Smith 标准形算法
% 纲要位置：计划纲要.md，第 1744 行。

% 纲要要点（供目录核对）：
%
% * 整数与一元多项式矩阵的 Smith 标准形及精确算法；保持矩阵等式与可逆性，区分精确商清除和余数降主元，以局部下降及有限主元阶段证明终止；正式证明整数证书的三个检查条件及原坐标的解、核、余核
% #### A.1.1 行列变换的记录
% #### A.1.2 主元约化与整除链
% #### A.1.3 标准形证书与线性方程
%

\section{整数与一元多项式矩阵的 Smith 标准形算法}
\label{b2:sec:A01}

\chapref{b2:ch:04}已经证明主理想整环上 Smith 标准形的存在性和唯一性。实际求解时，还要回答两个问题：允许哪些运算，怎样确认最终输出确实对应原来的矩阵？本节在整数环或有效给定的域上的一元多项式环中使用带余除法，记录每一步换基，使标准形同时给出原坐标中的解、核与余核。

\subsection{行列变换的记录}
\label{b2:subsec:A01:01}

设 \(R=\mathbb Z\) 或 \(R=k[t]\)，其中域 \(k\) 的加减乘除和相等判定均可精确执行。输入 \(A\in M_{m\times n}(R)\)，从
\[
B=A,\qquad P=I_m,\qquad Q=I_n
\]
出发，每一步都保持 \(B=PAQ\)、\(P\in\operatorname{GL}_m(R)\) 和 \(Q\in\operatorname{GL}_n(R)\)。作行变换 \(B\leftarrow EB\) 时同步令 \(P\leftarrow EP\)，因为 \(EB=(EP)AQ\)；作列变换 \(B\leftarrow BF\) 时同步令 \(Q\leftarrow QF\)，因为 \(BF=PA(QF)\)。这里 \(E,F\) 分别是相应尺寸的可逆初等矩阵，所以更新后的 \(P,Q\) 仍可逆。这些等式和可逆性是算法始终保持的不变量；行变换的矩阵累积在 \(P\) 的左侧，列变换的矩阵累积在 \(Q\) 的右侧。

\ProseParagraphBreak
这些规则看似只是在保存计算过程，却决定了最后能否把解送回原坐标。若 \(x=Qy\)，原方程 \(Ax=b\) 恰好变成 \(By=Pb\)。因此右端要随行变换改变，而未知量的坐标要用 \(Q\) 变回去，不能把两个换基矩阵混用。

\begin{example}\label{b2:exm:A-integer-certificate}
取整数矩阵
\[
A=\begin{pmatrix}4&6&2\\2&6&4\end{pmatrix}.
\]
先交换两行，再从第二行减去第一行的两倍，最后把第二行乘 \(-1\)，得到
\[
\begin{pmatrix}2&6&4\\0&6&6\end{pmatrix}.
\]
接着依次作 \(C_2\leftarrow C_2-3C_1\)、\(C_3\leftarrow C_3-2C_1\)、\(C_3\leftarrow C_3-C_2\)，其中最后一步的 \(C_2\) 指已经更新的列。结果及累积矩阵为
\begin{equation}\label{b2:eq:A-integer-certificate}
P=\begin{pmatrix}0&1\\-1&2\end{pmatrix},\qquad
Q=\begin{pmatrix}1&-3&1\\0&1&-1\\0&0&1\end{pmatrix},\qquad
PAQ=\begin{pmatrix}2&0&0\\0&6&0\end{pmatrix}.
\end{equation}
两换基矩阵的行列式均为一，且 \(2\mid6\)，所以这个等式是 Smith 标准形的完整证书。原矩阵的一阶子式最大公因子为二，三个二阶子式分别为 \(12,12,12\)，也与对角元的乘积相符。
\end{example}

\subsection{主元约化与整除链}
\label{b2:subsec:A01:02}

整数带余除法使用非负余数，多项式带余除法使用次数严格较小的余式。对非零主元 \(a\)，若同一列中有 \(b\)，写 \(b=qa+r\)，减去相应的 \(q\) 倍主元行。余数非零时交换这两行，新的主元是 \(r\)；余数为零时，\(q=b/a\) 是环中的精确商，行减法直接把该项清成零。处理同行元素则使用列变换。\cref{b2:alg:A-smith}中的 \(j\) 表示当前尚未处理的首个行列位置，每次只在编号至少为 \(j\) 的行、列中操作，保留此前已经完成的对角块。

\begin{algorithm}[htbp]
\caption{Smith 约化}
\label{b2:alg:A-smith}
\begin{algorithmsteps}{\(A\in M_{m\times n}(R)\)，其中 \(R=\mathbb Z\) 或 \(k[t]\)，\(k\) 为有效给定的域}{幺模矩阵 \(P,Q\) 及满足 \(PAQ=D\) 的 Smith 标准形 \(D\)}
\State 初始化 \(B=A\)、\(P=I_m\)、\(Q=I_n\)、\(j=1\)。行变换同步令 \(P\leftarrow EP\)，列变换同步令 \(Q\leftarrow QF\)，保持 \(B=PAQ\) 及 \(P,Q\) 可逆。
\While{\(j\le\min(m,n)\)}
\State 若右下剩余子块为零，结束；否则把其中一个非零项移到 \((j,j)\)。
\State 用带余除法检查第 \(j\) 列下方和第 \(j\) 行右方。发现非零余数就把它换成主元，并重新检查；全被主元整除时，先用精确商作行减法清除该列下方，再作列减法清除该行右方。
\State 若主元 \(a=B_{jj}\) 不整除右下子块某项 \(b=B_{ik}\)，把第 \(i\) 行加到第 \(j\) 行；对 \(b=qa+r\) 作 \(C_k\leftarrow C_k-qC_j\)，交换第 \(j,k\) 列，使非零余数 \(r\) 成为新主元，然后返回上一步。否则保留当前主元，令 \(j\leftarrow j+1\)。
\EndWhile
\State 把非零对角元归一化为正整数或首一多项式，仍同步更新换基矩阵；输出 \(P,Q,B\)。
\end{algorithmsteps}
\end{algorithm}

分块检查也确实会缩小主元。清除第 \(j\) 行、列后，若右下子块的 \((i,k)\) 项 \(b\) 不被 \(a=B_{jj}\) 整除，只看第 \(j,i\) 行和第 \(j,k\) 列，相关部分的变换为
\[
\begin{pmatrix}a&0\\0&b\end{pmatrix}
\xrightarrow{R_j\leftarrow R_j+R_i}
\begin{pmatrix}a&b\\0&b\end{pmatrix}
\xrightarrow{C_k\leftarrow C_k-qC_j}
\begin{pmatrix}a&r\\0&b\end{pmatrix},
\qquad b=qa+r,\quad r\ne0.
\]
交换这两列后，\(r\) 成为新主元。上式只画出有关的两行、两列；把第 \(i\) 行加到第 \(j\) 行，也会把第 \(i\) 行的其他条目带入第 \(j\) 行，所以必须返回行列检查，重新清除。这里暂时破坏的是当前行、列的零项，目的是取得更小主元；编号小于 \(j\) 的已完成块不受影响，因为剩余行在这些列上为零，剩余列在这些行上也为零。

\ProseParagraphBreak
固定 \(j\) 后，每次遇到非零余数都会严格缩小主元：整数情形为 \(0<r<|a|\)，多项式情形为 \(\deg r<\deg a\)。右下块的整除障碍经上述变换也造成同样的下降，因而主元只能更换有限次。在两次更换之间，只需检查有限个条目；若余数全部为零，先清列再清行的精确商操作也只有有限次，且清行不会重新产生已清掉的列下方条目。最终主元整除整个剩余子块，\(j\) 才增加。由于 \(j\le\min(m,n)\)，整个算法也在有限步后结束。

\ProseParagraphBreak
保留当前主元后，后续只在右下子块内作行列变换，新条目仍是旧条目的 \(R\) 线性组合，所以仍被这个主元整除。这保证后续得到的各个非零对角元接成整除链。最后的归一化也只乘单位：负整数主元所在的行乘 \(-1\)，多项式主元若首项系数为 \(c\ne0\)，所在的行乘 \(c^{-1}\)。因此正值或首一归一化不改变幺模等价，仍可同步保留完整的 \(P,Q\)。

\ProseParagraphBreak
这给出在 \(\mathbb Z\) 与 \(k[t]\) 上实现\cref{b2:thm:smith-normal-form}的算法，终止依据是可计算的绝对值或次数下降。一般主理想整环上的存在证明使用 Bézout 变换与主元理想的升链条件；要执行相应计算，还需具备\secref{b2:sec:A02}列出的有效运算。只做到对角化仍不够：例如 \(\operatorname{diag}(6,10)\) 不满足 Smith 条件，因为 \(6\nmid10\)；其一阶子式最大公因子为二，行列式为六十，\cref{b2:thm:smith-uniqueness}给出的归一化对角元应为 \(2,30\)。

\subsection{标准形证书与线性方程}
\label{b2:subsec:A01:03}

\begin{proposition}\label{b2:prop:A-smith-solving}
设 \(R\) 是主理想整环，\(A\in M_{m\times n}(R)\)，并已知
\[
PAQ=D,\qquad P\in\operatorname{GL}_m(R),\quad Q\in\operatorname{GL}_n(R),
\]
其中 \(D\) 是非零对角元为 \(d_1\mid\cdots\mid d_r\) 的 Smith 标准形。给定 \(b\in R^m\)，记 \(c=Pb\)。方程 \(Ax=b\) 有解当且仅当
\[
d_i\mid c_i\quad(1\le i\le r),\qquad c_i=0\quad(r<i\le m).
\]
有解时，全部解是 \(x=Qy\)，其中 \(y_i=c_i/d_i\) 对 \(i\le r\) 成立，其余 \(n-r\) 个坐标任取。特别地，若 \(e_i\) 是 \(R^n\) 的第 \(i\) 个标准基向量，则 \(\ker A\) 的一组基为 \(Qe_{r+1},\ldots,Qe_n\)。此外，有余核分解
\[
R^m/A(R^n)\cong\bigoplus_{i=1}^r R/(d_i)\ \oplus R^{m-r}.
\]
若 \(f_i\) 是 \(R^m\) 的第 \(i\) 个标准基向量，则第 \(i\) 个分量的坐标生成元在原余核中由 \(P^{-1}f_i+A(R^n)\) 代表：\(i\le r\) 时对应 \(R/(d_i)\)，\(i>r\) 时对应自由分量。
\end{proposition}
\begin{proof}
代入 \(x=Qy\) 并左乘 \(P\)，把原方程等价变为 \(Dy=c\)。前 \(r\) 行分别是 \(d_i y_i=c_i\)，其可解性就是整除条件；由于 \(d_i\ne0\) 且 \(R\) 为整环，可解时该坐标唯一。剩余行左端为零，正好给出所述零坐标条件；剩余列不受限制。\(Q\) 是同构，把这些零列对应的自由坐标轴送成原核的一组基。

由\cref{b2:prop:smith-cokernel}，余核同构把 \(v+A(R^n)\) 送到 \(Pv+D(R^n)\)；其逆把对角坐标 \(f_i\) 的类送到 \(P^{-1}f_i+A(R^n)\)。前 \(r\) 个坐标模掉 \((d_i)\)，余下零行的坐标则没有关系约束，所以得到所述循环分量与自由分量。若 \(d_i\) 是单位，对应循环分量为零。
\end{proof}

在\cref{b2:exm:A-integer-certificate}中取 \(b=(10,8)^{\mathsf T}\)，则 \(Pb=(8,6)^{\mathsf T}\)，所以
\[
x=\begin{pmatrix}1\\1\\0\end{pmatrix}
  +z\begin{pmatrix}1\\-1\\1\end{pmatrix},\qquad z\in\mathbb Z.
\]
直接乘原矩阵分别得到 \(b\) 和零，便同时核验了一个特解与齐次解。若把右端换成 \((8,8)^{\mathsf T}\)，则变换后的第二个坐标是八，不能被六整除，故没有整数解，即使在 \(\mathbb Q\) 上仍有解。

\ProseParagraphBreak
零列给出核的自由坐标，零行给出余核的自由坐标，两者在矩形矩阵中不能混同。若系数环有零因子，非零对角元的零化子也可能贡献核，不能只读取零列；\cref{b2:prop:A-diagonal-annihilators}将给出相应的一般公式。若只要抽象余核的同构类型，Smith 对角数据已经足够；恢复余核生成元的原坐标使用 \(P^{-1}\)，恢复未知量或核的原坐标则使用 \(Q\)。

\begin{proposition}\label{b2:prop:A-smith-certificate-validation}
设 \(m,n\ge1\) 为整数，\(A,D\in M_{m\times n}(\mathbb Z)\)，\(P\in M_m(\mathbb Z)\)、\(Q\in M_n(\mathbb Z)\)。若
\[
PAQ=D,\qquad \det P,\det Q\in\{1,-1\},
\]
且 \(D\) 是矩形对角矩阵，其非零对角元占据前 \(r\) 个位置，其中整数 \(r\) 与这些对角元满足
\[
0\le r\le\min(m,n),\qquad
d_i>0\ (1\le i\le r),\qquad d_1\mid\cdots\mid d_r,
\]
其余条目均为零，则 \((P,Q,D)\) 给出 \(A\) 的正归一化整数 Smith 标准形证书。\(D\) 由 \(A\) 唯一确定，而换基矩阵 \(P,Q\) 不要求唯一。验证这样的证书必须同时检查整数矩阵等式、整数上的可逆性以及归一化整除链。
\end{proposition}
\begin{proof}
整数矩阵的行列式为 \(\pm1\) 时，伴随矩阵公式给出的逆矩阵仍有整数条目，所以 \(P\in\operatorname{GL}_m(\mathbb Z)\)、\(Q\in\operatorname{GL}_n(\mathbb Z)\)。因此所给矩阵等式证明 \(A,D\) 在整数环上幺模等价。对角形、非零元排列和整除链满足\cref{b2:def:smith-normal-form}的全部条件，故 \(D\) 确为 Smith 标准形。\cref{b2:thm:smith-uniqueness}确定各非零对角元的相伴类，正值归一化再确定唯一代表。反过来，整数幺模矩阵的行列式必为 \(\pm1\)，而正归一化的 Smith 形必满足所写排列及整除条件；这些检查与计算标准形时采取了哪条变换路径无关。

例如取 \(A=D=\operatorname{diag}(2,3)\)、\(P=Q=I_2\)，矩阵等式与整数可逆性均成立，但 \(2\nmid3\)，所以这不是 Smith 证书。其条目最大公因子为一，行列式为六，由\cref{b2:thm:smith-uniqueness}，正归一化 Smith 形应为 \(\operatorname{diag}(1,6)\)。具体取
\[
P_0=\begin{pmatrix}1&1\\-3&-2\end{pmatrix},\qquad
Q_0=\begin{pmatrix}-1&-3\\1&2\end{pmatrix},\qquad
P_0\operatorname{diag}(2,3)Q_0=\operatorname{diag}(1,6).
\]
\(Q_0\) 的第一列使原对角矩阵的像为 \((-2,3)^{\mathsf T}\)，再由 \(P_0\) 的第一行读取 Bézout 等式 \(-2+3=1\)，取得单位主元；\(Q_0\) 的第二列使这一行的另一项为零。两换基矩阵的行列式都为一，所以这是合法的整数证书。另一方面，有理矩阵等式
\[
\operatorname{diag}(1/2,1/3)\operatorname{diag}(2,3)=I_2
\]
也不能给出整数 Smith 证书：左侧缩放矩阵根本不是整数矩阵，更不属于 \(\operatorname{GL}_2(\mathbb Z)\)。事实上，\(\operatorname{diag}(2,3)\) 的整数余核为 \(\mathbb Z/2\mathbb Z\oplus\mathbb Z/3\mathbb Z\ne0\)，而 \(I_2\) 的整数余核为零，故二者不可能在整数环上幺模等价。
\end{proof}

% !TeX root = ../../Book2.tex
% 纲要标题：### A.2 主理想整环的计算与一般环的障碍

% 纲要要点（供目录核对）：
% * 主理想整环上的计算和一般环上的障碍
% #### A.2.1 有效表示与终止条件
% #### A.2.2 非主理想造成的障碍
% #### A.2.3 系数增长与模素数检验

\section{主理想整环的计算与一般环的障碍}
\label{b2:sec:A02}

Smith 理论把主理想条件转化为结构结论，但计算还取决于环如何给出。这里先区分“存在所需元素”与“能够求出所需元素”，再用一个二元多项式矩阵说明：离开主理想整环后，失败的可能是标准形本身，而不只是现有算法太慢。

\subsection{有效表示与终止条件}
\label{b2:subsec:A02:01}

对给定主理想整环 \(R\)，执行第四章的 Bézout 约化至少需要以下操作：有限地表示输入元素，精确执行环运算和相等判定；判定 \(a\mid b\) 并在成立时求出商；对非同时为零的 \(a,b\)，求出
\[
g=ua+vb,\qquad a=ga',\qquad b=gb',
\]
其中 \(g\) 生成 \((a,b)\)；还需判定单位并求其逆。若要求唯一的元素输出，须进一步指定可执行的相伴类归一化。这里所说的有效计算条件是对环的表示及操作提出的要求，并不是主理想整环定义中额外隐含的一部分。

\ProseParagraphBreak
在这些操作已给出的前提下，矩阵
\[
E=\begin{pmatrix}u&v\\-b'&a'\end{pmatrix}
\]
满足 \(\det E=ua'+vb'=1\)，并把 \((a,b)^{\mathsf T}\) 化为 \((g,0)^{\mathsf T}\)。若 \(a\nmid b\)，主元理想从 \((a)\) 严格扩大为 \((g)\)。\cref{b2:lem:pid-ascending-chain}已经证明主理想整环中的理想升链最终稳定，因此反复使用这样的约化也会终止，不必预先找到 Euclid 函数。这个论证保证给定输入在有限步后结束，却没有给出随输入长度变化的统一运行时间界。

\ProseParagraphBreak
这里说域 \(k\) 有效给定，是指它的元素有有限表示，且加、减、乘、非零元的除法以及相等和零判定都能在有限步内执行。整数和这样的域上的一元多项式环满足上述计算条件。例如，多项式扩展 Euclid 算法不仅输出最大公因式，还逐步保存 Bézout 系数；首一归一化只需除以一个非零域元素。这一过程及次数估计可参阅 Shoup 的《A Computational Introduction to Number Theory and Algebra》\cite[Section 17.3, Theorems 17.2, 17.4, pp. 469--472]{Shoup2008NumberTheoryAlgebra}。域的有效性不能省略：给出一个任意实数的近似值，并不同时给出了精确的零判定。

\subsection{非主理想造成的障碍}
\label{b2:subsec:A02:02}

\begin{proposition}\label{b2:prop:A-nonprincipal-obstruction}
设 \(k\) 为域，\(R=k[x,y]\)。行矩阵 \(A=(x\ y)\) 不与任何 \((d\ 0)\) 在 \(R\) 上幺模等价。
\end{proposition}
\begin{proof}
可逆行列变换保持所有条目生成的理想。事实上，变换后每个条目是原条目的 \(R\) 线性组合，给出一个包含关系；用逆变换得到反向包含。因此若这样的等价存在，便有 \((x,y)=(d)\)。

由此 \(d\) 同时整除 \(x,y\)。写 \(x=dh\)，并把等式看作 \(k[x][y]\) 中关于 \(y\) 的多项式等式。由于 \(x\ne0\) 且非零多项式的 \(y\) 次数在乘法下相加，有 \(\deg_y d=0\)，即 \(d\in k[x]\)。再写 \(y=dg\)，比较 \(y\) 的一次项系数，得到 \(d\) 在 \(k[x]\) 中整除 \(1\)，故 \(d\) 是单位。这会给出 \((x,y)=R\)。但代入 \(x=y=0\)，每个 \(xf+yg\) 都取零值，不可能等于一，矛盾。
\end{proof}

这个例子在域 \(k(x,y)\) 上只有秩一，仍不能在原环上化成一个非零对角元。秩丢失了非主理想 \((x,y)\) 的信息。一般交换环上的条目理想和高阶子式理想仍有意义，但不能预设每个这样的理想都有单个生成元；第三卷第12.3节有关模展示与 Fitting 理想的讨论将继续利用这些不变量。

\ProseParagraphBreak
有零因子时还要另查核的计算。例如在 \(R=\mathbb Z/6\mathbb Z\) 中，对角方程 \(2z=0\) 有两个解 \(z=0,3\)。因而即使矩阵已经对角化，也不能照用\cref{b2:prop:A-smith-solving}中“非零对角元的对应齐次坐标为零”这一步；该步明确使用了整环的非零消去律。

\subsection{系数增长与模素数检验}
\label{b2:subsec:A02:03}

整数运算的代价不仅取决于运算次数，还取决于中间整数的位数。行列相减虽不会引入分母，倍数的累积仍可能使换基矩阵的条目很大。对有理数先清除分母，可以把零空间或线性方程化为整数计算；若随后研究整数余核，却必须保留这一缩放的含义，因为乘一个非单位会改变整数关系模。

\begin{proposition}\label{b2:prop:A-modular-rank}
设 \(A\in M_{m\times n}(\mathbb Z)\)，\(p\) 是素数，\(\overline A\in M_{m\times n}(\mathbb F_p)\) 是将 \(A\) 的条目逐一模 \(p\) 所得的矩阵，则
\[
\operatorname{rank}_{\mathbb F_p}\overline A\le
\operatorname{rank}_{\mathbb Q}A.
\]
若某个 \(r\) 阶子式在模 \(p\) 后非零，则原矩阵的有理数秩至少为 \(r\)。若同时给出了 \(A=UV\)，其中 \(U\) 为 \(m\times r\) 整数矩阵、\(V\) 为 \(r\times n\) 整数矩阵，那么原矩阵的有理数秩恰为 \(r\)。
\end{proposition}
\begin{proof}
\(\overline A\) 的非零 \(r\) 阶子式来自一个不被 \(p\) 整除的整数子式，所以该整数非零，原矩阵秩至少为 \(r\)。取 \(r\) 为模 \(p\) 的秩即得不等式。另一方面，若 \(A=UV\)，它在 \(\mathbb Q\) 上的像包含于 \(U\) 的至多 \(r\) 维的列空间，所以秩至多为 \(r\)，与前面的下界合并得到等号。
\end{proof}

例如 \(\operatorname{diag}(1,6)\) 在 \(\mathbb Q\) 上秩为二，模二或模三后秩降为一，模五后秩仍为二。模 \(p\) 为零的整数子式未必在整数中为零；模 \(p\) 非零的子式则能证明对应的整数子式非零。

\begin{proposition}\label{b2:prop:A-bounded-congruence}
设 \(B\ge0\) 为实数，\(s\ge1\) 为整数，\(z\in\mathbb Z\)、\(|z|\le B\)，并设两两互素的正整数 \(q_1,\ldots,q_s\) 满足
\[
z\equiv0\pmod {q_i}\quad(1\le i\le s),\qquad
q_1\cdots q_s>B.
\]
则 \(z=0\)。若 \(z,z'\in[-B,B]\cap\mathbb Z\) 在这些模数下逐一同余，且 \(q_1\cdots q_s>2B\)，则 \(z=z'\)。
\end{proposition}
\begin{proof}
两两互素性说明乘积 \(M=q_1\cdots q_s\) 整除 \(z\)。非零的 \(M\) 倍数绝对值至少为 \(M\)，与 \(|z|\le B<M\) 矛盾。第二个结论对 \(z-z'\) 使用第一个结论，因为 \(|z-z'|\le2B<M\)。
\end{proof}

这使模运算也能提供确定性的等式证书。检验整数矩阵的等式 \(PAQ=D\) 时，可以直接从输入条目取得上界，无须先算出整个乘积。设 \(m,n\geq1\)，\(A,D\in M_{m\times n}(\mathbb Z)\)、\(P\in M_m(\mathbb Z)\)、\(Q\in M_n(\mathbb Z)\)，并记 \(\lVert M\rVert_{\max}=\max_{i,j}\lvert M_{ij}\rvert\)。每个乘积条目是 \(mn\) 项的和，因而
\[
\begin{aligned}
B&\coloneqq mn\lVert P\rVert_{\max}\cdot\lVert A\rVert_{\max}
\cdot\lVert Q\rVert_{\max}+\lVert D\rVert_{\max},\\
\lVert PAQ-D\rVert_{\max}&\leq B.
\end{aligned}
\]
这个粗界只需计算四个输入矩阵的最大条目绝对值。选取乘积大于 \(B\) 的两两互素模数，并在每个模数下验证 \(PAQ=D\)，\cref{b2:prop:A-bounded-congruence}就逐条目保证整数等式成立。

\ProseParagraphBreak
没有上界时，即使多次模检验都给出零，也不能断定整数等式成立：模数的乘积本身非零，却在每个模数下都为零。用中国剩余定理恢复有界整数的矩阵计算方法，见 Shoup 的上述著作\cite[Section 4.4, pp. 84--86]{Shoup2008NumberTheoryAlgebra}；这里仅使用有界整数的唯一性，不对本附录的 Smith 算法声称复杂度最优。

% !TeX root = ../../Book2.tex
% 纲要标题：### A.3 精确代数结论与近似输入的区别
% 纲要位置：计划纲要.md，第 1764 行。

% 纲要要点（供目录核对）：
%
% * 精确结构分类与近似输入所能支持的结论；区别输入区间不确定性与固定输入的舍入误差，有限二进制浮点数本身与其所代表的目标实数分别说明
% <!-- 短节：不设小节 -->
%

\section{精确代数结论与近似输入的区别}
\label{b2:sec:A03}

精确代数计算处理指定环或域中的恒等式、整除与相等。与近似计算比较时，须区分输入本身的不确定性和算法的舍入误差：只知道某个参数所在的区间，与已经固定一个输入、但运算过程中发生舍入，是两个不同问题。本节主要讨论前者。例如对实参数矩阵
\[
A_\varepsilon=\begin{pmatrix}1&1\\1&1+\varepsilon\end{pmatrix},
\]
精确行列式为 \(\varepsilon\)。因此 \(\varepsilon=0\) 时秩为一，而任意非零 \(\varepsilon\) 都给出秩二。如果输入只知道 \(|\varepsilon|\le10^{-12}\)，其中既包含零，也包含非零数；增加后续计算的精度不能消除输入本身留下的这一不确定性。

\ProseParagraphBreak
同样，矩阵
\[
N=\begin{pmatrix}0&1\\0&0\end{pmatrix},\qquad
N_\varepsilon=\begin{pmatrix}0&1\\\varepsilon&0\end{pmatrix}
\]
满足 \(N_\varepsilon\to N\)（\(\varepsilon\to0^+\)）。\(N\) 的最小多项式是 \(t^2\)，有一个二阶 Jordan 块；当实数 \(\varepsilon>0\) 时，\(N_\varepsilon\) 有两个不同的实特征值 \(\sqrt\varepsilon,-\sqrt\varepsilon\)，因此在实数域上可对角化，其 Jordan 形由这两个一阶块组成。Jordan 型中的块大小与数量是精确相似分类的离散数据；这个例子表明，任意小的条目扰动都可能改变它们。因而不能凭近似特征值是否相等来判断精确块结构，而须使用第四章的不变因子或核维数论证。

\ProseParagraphBreak
这也不表示精确算法能够从测量值恢复唯一的真实模型。一个有限的二进制浮点数本身是确定的数，可以精确解释为某个整数乘以二的整数次幂；近似性来自用这个数代表目标实数，浮点运算还可能引入新的舍入。把十进制字符串 \(0.1\) 指定为有理数 \(1/10\)，是在选择一个精确输入；通常的有限二进制浮点存储不能恰好表示 \(1/10\)，保存的是另一个二进有理数。若把该字符串解释为区间 \([0.0999,0.1001]\) 中的未知实数，输入又不再是单个确定值。显示相同的小数不保证这几种解释指定了同一个数学对象，题目应先明确采用哪一种。

\ProseParagraphBreak
在本书的代数算例中，整数以整数运算处理，有理数以分子分母处理，有限域元素按明确给出的模关系处理。输出应能通过精确乘法、余式为零、换基行列式在工作环中为单位等条件复核。数值计算则可以用于发现候选结构、估计解或提示某个矩阵接近降秩；若要把这些观察提升为精确命题，还须另给足够的代数证据。本节不讨论误差传播或数值稳定性的算法分析，这些问题属于数值线性代数的范围。

% !TeX root = ../../Book2.tex
% 纲要标题：### A.4 有限域与有理数上的精确算例
% 纲要位置：计划纲要.md，第 1770 行。

% 纲要要点（供目录核对）：
%
% * 有限域和有理数上的可复核例题；形式多项式与有限点函数、目标余核与来源坐标、环生成数与域维数、分式域可逆与原环满射分别核对；有理相似只须非零行列式，先证明双侧逆再给逆公式；正式证明全部逐点秩仍不能恢复重数
% #### A.4.1 有限域上的多项式矩阵
% #### A.4.2 有理数上的循环基与相似证书
%

\section{有限域与有理数上的精确算例}
\label{b2:sec:A04}

下面的两个算例分别计算多项式矩阵的余核和有理矩阵的相似类。前者使用左右两个独立的换基，后者要求同一线性空间在定义域与值域采用同一组新基。把这两种等价关系混在一起，即使矩阵乘法都算对了，也会回答另一个问题。

\subsection{有限域上的多项式矩阵}
\label{b2:subsec:A04:01}

取 \(R=\mathbb F_5[t]\)，考虑
\[
B=\begin{pmatrix}t^2+1&t\\t&t^2\end{pmatrix}.
\]
所有系数在 \(\mathbb F_5\) 中运算，但 \(t\) 仍是形式多项式变量，不把它代成域中某个数。\(\mathbb F_5[t]\) 与有限集合 \(\mathbb F_5\) 上的函数环不是同一个对象：非零多项式 \(t^5-t\) 在每个 \(a\in\mathbb F_5\) 处的值都为零，所以逐点求值会把不同多项式变成同一个函数。下面的变换均在多项式环中进行。先作 \(C_1\leftarrow C_1-tC_2\)，得到
\[
\begin{pmatrix}1&t\\t-t^3&t^2\end{pmatrix}.
\]
再作 \(R_2\leftarrow R_2+(t^3-t)R_1\)，清除第一列下项，最后作 \(C_2\leftarrow C_2-tC_1\)。这给出
\begin{equation}\label{b2:eq:A-polynomial-certificate}
P=\begin{pmatrix}1&0\\t^3-t&1\end{pmatrix},\qquad
Q=\begin{pmatrix}1&-t\\-t&1+t^2\end{pmatrix},\qquad
D=PBQ=\begin{pmatrix}1&0\\0&t^4\end{pmatrix}.
\end{equation}
这里 \(\det P=\det Q=1\)，所以换基始终在 \(\mathbb F_5[t]\) 上可逆。独立检查中，一阶子式生成单位理想，因为 \((t^2+1)-t\cdot t=1\)；行列式为 \(t^4\)。由\cref{b2:thm:smith-uniqueness}，不变因子确为 \(1,t^4\)。

\ProseParagraphBreak
由于 \(Q\) 可逆，\(PB(R^2)=D(R^2)\)，所以 \(P\) 诱导同构
\[
\operatorname{coker}B\xrightarrow{\sim}\operatorname{coker}D,
\qquad x+B(R^2)\longmapsto Px+D(R^2).
\]
对 \(D=\operatorname{diag}(1,t^4)\)，取第二坐标模去 \(t^4\) 又给出 \(\operatorname{coker}D\cong R/(t^4)\)，与\cref{b2:prop:smith-cokernel}一致。两个同构的复合在原自由模上由映射
\[
R^2\longrightarrow R/(t^4),\qquad
\binom{u}{v}\longmapsto (t^3-t)u+v+(t^4)
\]
给出；其核恰是 \(B(R^2)\)。标准余核中的生成元 \(e_2+D(R^2)\) 经逆同构回到
\[
P^{-1}e_2+B(R^2).
\]
本例恰有 \(P^{-1}e_2=e_2\)，所以原坐标中 \(e_2\) 的陪集就生成这个循环模；一般恢复生成元时须实际计算 \(P^{-1}e_i\)，不能略掉这一换基。\(Q\) 则改变来源自由模中的关系坐标，其可逆性使关系列表的重组不改变 \(B(R^2)\)。商模作为 \(\mathbb F_5\) 向量空间的一组基为
\[
\overline e_2,\quad t\overline e_2,\quad
t^2\overline e_2,\quad t^3\overline e_2.
\]
乘 \(t\) 的算子在这组基下是 \(C(t^4)\)，其最小多项式为 \(t^4\)。作为 \(R\) 模，只需一个生成元 \(\overline e_2\)，因为 \(R\) 系数已允许乘上 \(t\) 的任意多项式；作为 \(\mathbb F_5\) 向量空间，却需要上述四个基元素。这两个生成数使用的系数环不同。

\ProseParagraphBreak
若只计算 \(B\) 在分式域 \(\mathbb F_5(t)\) 上的秩，会得到二：行列式 \(\det B=t^4\) 在该域中是单位，所以矩阵在那里可逆。但 \(t^4\) 不是 \(R=\mathbb F_5[t]\) 中的单位；非零常数才是这个多项式环的单位。因此不能据分式域上的可逆性推出原映射在 \(R\) 上满射，原余核仍是非零的 \(R/(t^4)\)。

\subsection{有理数上的循环基与相似证书}
\label{b2:subsec:A04:02}

在 \(\mathbb Q^3\) 上取算子矩阵
\[
T=\begin{pmatrix}
\tfrac12&-\tfrac14&3\\
1&-\tfrac12&2\\
0&1&0
\end{pmatrix},\qquad v=\begin{pmatrix}1\\0\\0\end{pmatrix}.
\]
逐次乘法给出
\[
Tv=\begin{pmatrix}\tfrac12\\1\\0\end{pmatrix},\qquad
T^2v=\begin{pmatrix}0\\0\\1\end{pmatrix},\qquad
T^3v=\begin{pmatrix}3\\2\\0\end{pmatrix}=2v+2Tv.
\]
前三个向量组成列矩阵
\[
S=(v\ Tv\ T^2v)=
\begin{pmatrix}1&\tfrac12&0\\0&1&0\\0&0&1\end{pmatrix},\qquad \det S=1.
\]
因此它们是一组循环基。在 \(\mathbb Q\) 上，合法换基只要求 \(\det S\ne0\)，本例行列式等于一不是相似变换所必需的条件；非零有理数，如二或 \(1/3\)，也能作为换基矩阵的行列式。这与整数 Smith 变换要求行列式为 \(\pm1\) 不同。记 \(f(t)=t^3-2t-2\)，则其友矩阵为
\[
C(f)=\begin{pmatrix}0&0&2\\1&0&2\\0&1&0\end{pmatrix}.
\]
三列所表示的等式合并为如下相似证书：
\begin{equation}\label{b2:eq:A-rational-similarity}
TS=SC(f),\qquad \det S=1,
\qquad S^{-1}TS=C(f).
\end{equation}
由相似证书，特征多项式为 \(\chi_T(t)=f(t)\)。最小多项式也能直接核验：次数至多为二的非零多项式不可能零化 \(v\)，否则 \(v,Tv,T^2v\) 线性相关；而 \(f(T)v=0\)，与 \(T\) 交换后得到 \(f(T)T^jv=0\)，所以 \(f(T)\) 零化整个空间。最小多项式因而次数为三，并整除首一三次多项式 \(f\)，故 \(\mu_T(t)=\chi_T(t)=f(t)\)，有理标准形只有这一个块。

\ProseParagraphBreak
还可以用同一证书求逆。先将 \(T^3-2T-2I=0\) 改写为
\[
T(T^2-2I)=2I.
\]
\(T^2-2I\) 是 \(T\) 的多项式，与 \(T\) 交换，所以另一顺序的乘积也为 \(2I\)。因此 \(\dfrac12(T^2-2I)\) 已被证明是 \(T\) 的双侧逆，才得到
\begin{equation}\label{b2:eq:A-rational-inverse}
T^{-1}=\tfrac12(T^2-2I).
\end{equation}
全程只使用有理数运算，没有求近似特征值，也不必先把 \(f\) 的根写出来。

\ProseParagraphBreak
\eqref{b2:eq:A-polynomial-certificate}中的 \(P,Q\) 是两个自由模的独立换基，保存矩阵给出的余核；\eqref{b2:eq:A-rational-similarity}中的 \(S^{-1}TS\) 是同一空间的换基，保存算子的相似类。若把后一种换基写成左右乘法，左右矩阵必须互逆，即 \(P=S^{-1},Q=S\)。独立的行列操作一般不保持特征多项式。判断余核同构或线性方程的可解性时使用前一种证书，判断算子相似性时使用后一种证书。

\ProseParagraphBreak
逐点求值能检查某些性质，却不能恢复多项式矩阵的全部不变因子。即使在所有扩域中知道每个点的秩，零点处的重数仍可能没有被记录。

\begin{proposition}\label{b2:prop:A-pointwise-ranks-lose-multiplicity}
设 \(k\) 为任意域，\(r,s\) 为不同的正整数，\(R=k[t]\)，并令
\[
B_r=\operatorname{diag}(1,t^r),\qquad
B_s=\operatorname{diag}(1,t^s)
\]
表示从 \(R^2\) 到 \(R^2\) 的 \(R\) 线性映射。对每个扩域 \(K/k\) 和每个 \(a\in K\)，特殊化矩阵 \(B_r(a)\)、\(B_s(a)\) 的秩相同；但
\[
\operatorname{coker}B_r\cong k[t]/(t^r),\qquad
\operatorname{coker}B_s\cong k[t]/(t^s)
\]
的 \(k\) 维数分别为 \(r,s\)，所以这两个余核作为 \(R\) 模不同构。特别地，当 \(k=\mathbb F_5\)、\(r=4,s=8\) 时，两矩阵在每个 \(a\in\mathbb F_5\) 处的取值甚至完全相同。
\end{proposition}
\begin{proof}
对任意 \(K/k\) 和 \(a\in K\)，若 \(a=0\)，两个矩阵都为 \(\operatorname{diag}(1,0)\)，秩为一；若 \(a\ne0\)，正整数次幂 \(a^r,a^s\) 均非零，两个矩阵的秩都为二。因此全部逐点秩相同。

在 \(R\) 上，第一坐标的像为整个 \(R\)，第二坐标的像分别为 \(t^rR,t^sR\)。取第二坐标的商，便得到所列余核同构。由首一多项式的带余除法，\(R/(t^r)\) 的 \(k\) 基为 \(1,t,\ldots,t^{r-1}\) 的商类，而 \(R/(t^s)\) 有 \(s\) 个相应基元素。任何 \(R\) 模同构也必须保持 \(k\) 标量作用，因而给出 \(k\) 线性同构；维数不同排除这种可能。

最后在 \(\mathbb F_5\) 中，\(a=0\) 时 \(a^4=a^8=0\)；乘法群 \(\mathbb F_5^\times\) 的阶为四，所以 \(a\ne0\) 时 \(a^4=1\)，进而 \(a^8=1\)。故 \(B_4(a)=B_8(a)\) 对每个 \(a\in\mathbb F_5\) 成立，而前段已证明这两个余核不同构。这同时区分了有限点上的矩阵函数和多项式矩阵本身。
\end{proof}


\begin{chaptersummary}
在指定的工作环上，矩阵等式 \(PAQ=D\)、换基矩阵的可逆性以及对角元的整除链和归一化共同构成 Smith 标准形证书。余核中目标向量的坐标由 \(P\) 给出，标准余核生成元由 \(P^{-1}\) 换回；方程右端变成 \(Pb\)，解与核则由 \(Q\) 换回。算法的终止理由来自 Euclid 量的下降或主元理想的升链条件；要真正执行它，还须能够计算整除关系、商和 Bézout 系数。

\ProseParagraphBreak
一般环上的非主理想可能阻止对角化，零因子则使非零对角元具有非平凡核。模素数计算可以快速提供秩的下界，有了明确的整数上界还可以证明精确等式；没有这些条件时，若干数值或同余观察不足以恢复全部结构。余核分类使用独立的左右换基，算子分类使用相似变换，这一区别应当在动手计算之前明确。
\end{chaptersummary}

\BookExercises

\begin{exercise}\label{b2:ex:A-wrong-diagonal}
令
\[
A=\operatorname{diag}(1,6),\qquad P=\operatorname{diag}(2,1),
\qquad Q=I_2,\qquad D=\operatorname{diag}(2,6).
\]
核验 \(PAQ=D\) 以及 \(D\) 的正数整除链，并证明 \(P\) 模每个奇素数都可逆。这些检查能否给出整数 Smith 标准形证书？比较 \(\operatorname{coker}A\) 与 \(\operatorname{coker}D\)，指出缺少的条件。
\end{exercise}
\begin{solution}
直接相乘得 \(PAQ=D\)，且 \(2\mid6\)。对任意奇素数 \(p\)，\(\det\bar P=\bar2\ne0\)，故 \(\bar P\) 在 \(\mathbb F_p\) 上可逆。但 \(\det P=2\) 不是整数环的单位，\(P\notin\operatorname{GL}_2(\mathbb Z)\)。因此等式和对角条件都已满足，缺少的是整数可逆性。

事实上
\[
\operatorname{coker}A\cong\mathbb Z/6\mathbb Z,
\qquad
\operatorname{coker}D\cong\mathbb Z/2\mathbb Z\oplus\mathbb Z/6\mathbb Z,
\]
两者分别有六和十二个元素，并不同构。即使检查了所有奇素数下的可逆性，也漏掉了行列式中的因子二；模二检验会发现这处障碍。
\end{solution}

\begin{exercise}\label{b2:ex:A-original-coordinates}
对\cref{b2:exm:A-integer-certificate}中的矩阵 \(A\)，给出方程 \(Ax=(a,b)^{\mathsf T}\) 存在整数解的充要条件和全部整数解。再在原余核中找出阶分别为二和六的两个生成元，使其给出循环直和分解。
\end{exercise}
\begin{solution}
\eqref{b2:eq:A-integer-certificate}给出 \(P(a,b)^{\mathsf T}=(b,2b-a)^{\mathsf T}\)。由\cref{b2:prop:A-smith-solving}，充要条件为
\[
2\mid b,\qquad6\mid(2b-a).
\]
成立时，先取 \(y=\bigl(b/2,(2b-a)/6,z\bigr)^{\mathsf T}\)，再乘 \(Q\)，得到
\[
x=\begin{pmatrix}(a-b)/2\\(2b-a)/6\\0\end{pmatrix}
  +z\begin{pmatrix}1\\-1\\1\end{pmatrix},\qquad z\in\mathbb Z.
\]
上述整除条件保证显示的分数都是整数。

又 \(P^{-1}=\begin{pmatrix}2&-1\\1&0\end{pmatrix}\)，所以 \((2,1)^{\mathsf T}\) 和 \((-1,0)^{\mathsf T}\) 的陪集分别对应对角余核的两个标准生成元。由\cref{b2:prop:smith-cokernel}的具体同构，它们的阶恰为二与六，并且给出 \(\mathbb Z/2\mathbb Z\oplus\mathbb Z/6\mathbb Z\) 的直和。
\end{solution}

\begin{exercise}\label{b2:ex:A-diagonal-invariant-factors}
求 \(\operatorname{diag}(12,18,30)\) 的正整数 Smith 因子，说明为什么不能只把原对角元按大小排序。求其整数余核的元素个数、指数以及最少生成元个数。
\end{exercise}
\begin{solution}
一阶子式的最大公因子为六；非零二阶子式是 \(216,360,540\)，最大公因子为三十六；行列式为 \(6480\)。由\cref{b2:thm:smith-uniqueness}，三个不变因子是
\[
6,\qquad36/6=6,\qquad6480/36=180.
\]
大小排序不能代替整除链，例如 \(12<18\) 却有 \(12\nmid18\)。

余核为 \(G=\mathbb Z/6\mathbb Z\oplus\mathbb Z/6\mathbb Z\oplus\mathbb Z/180\mathbb Z\)，故元素个数为 \(6480\)，指数为 \(180\)。三个标准生成元足够，而 \(G/2G\) 是三维的 \(\mathbb F_2\) 向量空间，任何生成集的像也必须生成它，所以少于三个不够。
\end{solution}

\begin{exercise}\label{b2:ex:A-polynomial-bezout}
令 \(a=t^2+t+1\)、\(b=t^2-t+1\)。在 \(\mathbb F_5[t]\) 中求出 Bézout 系数和行列式为一的矩阵，把 \((a,b)^{\mathsf T}\) 化为 \((1,0)^{\mathsf T}\)。再在 \(\mathbb Z[t]\) 中求条目理想 \((a,b)\)，证明其商环同构于四元域，并说明为什么原整数多项式列向量不能通过可逆行变换化为 \((1,0)^{\mathsf T}\)。
\end{exercise}
\begin{solution}
展开得到 \((1-t)a+(t+1)b=2\)。在 \(\mathbb F_5\) 中二的逆是三，所以令 \(u=3(1-t)\)、\(v=3(t+1)\)，就有 \(ua+vb=1\)，从而
\[
E=\begin{pmatrix}u&v\\-b&a\end{pmatrix},\qquad
E\binom ab=\binom10,
\qquad\det E=ua+vb=1.
\]

在 \(\mathbb Z[t]\) 中，前一恒等式给出 \(2\in(a,b)\)，而 \(b=a-2t\)，所以 \((a,b)=(2,t^2+t+1)\)。于是
\[
\mathbb Z[t]/(a,b)\cong\mathbb F_2[t]/(t^2+t+1).
\]
右边的二次多项式在 \(\mathbb F_2\) 中没有根，故不可约，商环是四元域。特别地，\((a,b)\) 是真理想；可逆行变换保持条目理想，而 \((1,0)^{\mathsf T}\) 的条目生成整个环，所以不可能有这样的整数多项式换基。上述 \(u,v\) 若原样作为整数多项式使用，行列式是六，不能充当可逆变换。
\end{solution}

\begin{exercise}\label{b2:ex:A-integral-polynomial-obstruction}
证明 \(\mathbb Z[t]\) 上的行矩阵 \((2\ t)\) 不能通过幺模变换化成 \((d\ 0)\)。比较它在 \(\mathbb Q[t]\) 上的情形。
\end{exercise}
\begin{solution}
若能化成所述形状，条目理想不变性给出 \((2,t)=(d)\)。由 \(d\mid2\)，次数相加说明 \(d\) 是整数常数；再由 \(d\mid t\)，比较 \(t\) 的系数得 \(d\mid1\)，故 \(d=\pm1\)。但映射 \(\mathbb Z[t]\to\mathbb F_2\)、\(f\mapsto f(0)\bmod2\) 把 \((2,t)\) 送为零，故该理想不含一，矛盾。

在 \(\mathbb Q[t]\) 中二为单位，取
\[
Q=\begin{pmatrix}1/2&-t/2\\0&1\end{pmatrix}
\]
即有 \((2\ t)Q=(1\ 0)\)，且 \(\det Q=1/2\) 是该环的单位。改变系数环会改变允许的换基与余核。
\end{solution}

\begin{exercise}\label{b2:ex:A-zero-divisor-kernel}
在 \(R=\mathbb Z/12\mathbb Z\) 上令
\[
A=\begin{pmatrix}4&0&4\\0&6&6\end{pmatrix},\qquad
Q=\begin{pmatrix}1&0&-1\\0&1&-1\\0&0&1\end{pmatrix}.
\]
核验 \(AQ\) 为矩形对角矩阵，计算 \(A:R^3\to R^2\) 的核、像与余核，并证明核不是自由 \(R\) 模。最后给出 \(Ax=(4,6)^{\mathsf T}\) 的全部解。
\end{exercise}
\begin{solution}
\(\det Q=1\)，直接相乘得 \(AQ=\begin{pmatrix}4&0&0\\0&6&0\end{pmatrix}\)。令 \(x=Qy=(y_1-y_3,y_2-y_3,y_3)^{\mathsf T}\)。\(Ax=0\) 等价于 \(4y_1=0\)、\(6y_2=0\)，所以核中的向量恰为
\[
\begin{pmatrix}3a-c\\2b-c\\c\end{pmatrix},
\qquad 0\le a<4,\quad0\le b<6,\quad0\le c<12,
\]
各坐标在 \(R\) 中读取。由第三坐标恢复 \(c\)，再由 \(x_1+c=3a\)、\(x_2+c=2b\) 分别恢复 \(a\bmod4\)、\(b\bmod6\)，所以这些表示互不重复。映射 \(r\mapsto3r\) 的像是 \(\operatorname{Ann}_R(4)\)，核为 \((4)\)；映射 \(r\mapsto2r\) 的像是 \(\operatorname{Ann}_R(6)\)，核为 \((6)\)。由\cref{b2:thm:module-first-isomorphism}，分别得到 \(R/(4)\cong\operatorname{Ann}_R(4)\) 和 \(R/(6)\cong\operatorname{Ann}_R(6)\)，于是
\[
\ker A\cong R/(4)\oplus R/(6)\oplus R,
\qquad
\operatorname{im}A=(4)\oplus(6),
\qquad
\operatorname{coker}A\cong R/(4)\oplus R/(6).
\]
因为 \(R/(4)\)、\(R/(6)\) 分别有四、六个元素，而理想 \((4)\)、\((6)\) 分别有三、两个元素，三者的元素数分别为 \(4\cdot6\cdot12=288\)、\(3\cdot2=6\)、\(4\cdot6=24\)。有限秩自由 \(R\) 模的元素个数为 \(12^d\)，而 \(288\) 不是十二的整数次幂；无限秩自由模则有无限多个元素。因此核不是自由模。

对指定右端，\(y_1=y_2=1\)、\(y_3=0\) 是一组解。加上全部齐次解，再乘 \(Q\)，得到全部原坐标解
\[
x=\begin{pmatrix}1+3a-c\\1+2b-c\\c\end{pmatrix},
\qquad0\le a<4,\quad0\le b<6,\quad0\le c<12.
\]
此处不仅零列给出自由坐标，两个非零对角元也各自给出非唯一坐标。
\end{solution}

\begin{exercise}\label{b2:ex:A-modular-rank-certificate}
对整数矩阵
\[
A=\begin{pmatrix}2&4&6\\4&8&12\\6&12&20\end{pmatrix},
\]
给出它在 \(\mathbb Q\) 上秩为二的证书，证书须由一个模五的非零子式及一个 \(3\times2\) 与 \(2\times3\) 矩阵的乘积组成。模二的计算在这里能提供什么信息？
\end{exercise}
\begin{solution}
第一、三行与第一、三列组成的子式为 \(40-36=4\)，模五非零，所以\cref{b2:prop:A-modular-rank}给出有理数秩至少为二。另一方面，直接相乘有
\[
A=\begin{pmatrix}1&0\\2&0\\3&1\end{pmatrix}
  \begin{pmatrix}2&4&6\\0&0&2\end{pmatrix},
\]
所以秩至多为二，合并得到等于二。模二后矩阵为零，只给出无信息的下界零，不能据此断定原矩阵的秩为零。
\end{solution}

\begin{exercise}\label{b2:ex:A-modular-identity-bound}
设 \(P,C\) 是 \(m\) 阶整数矩阵，\(m\ge1\)，所有条目的绝对值至多为整数 \(H\ge1\)。给出 \(PC-I_m\) 每个条目的统一界。对两两互素整数 \(q_1,\ldots,q_s\ge2\)，已知 \(PC\equiv I_m\pmod{q_i}\)；给出模数乘积的一个充分条件，使这些同余证明 \(P\in\operatorname{GL}_m(\mathbb Z)\) 且 \(C=P^{-1}\)。若只知道 \(P\) 模这些模数可逆，而没有整数候选逆和大小界，结论是否仍成立？
\end{exercise}
\begin{solution}
每个 \((PC)_{ij}\) 是 \(m\) 个绝对值至多 \(H^2\) 的整数之和，因此
\[
\bigl|(PC-I_m)_{ij}\bigr|\le mH^2+1.
\]
若 \(M=\prod_iq_i>mH^2+1\)，则对每个条目应用\cref{b2:prop:A-bounded-congruence}，得到整数等式 \(PC=I_m\)。取行列式得
\[
\det P\det C=1.
\]
两因子都是整数，故均为 \(\pm1\)，\(P\) 在整数环上可逆；再左乘其逆得 \(C=P^{-1}\)，也就有 \(CP=I_m\)。这里同余和大小界不只验证了某个 Smith 乘积等式，而是验证了整数候选逆。

单有有限多个模数下的可逆性仍不足够。取 \(P=\operatorname{diag}(M+1,1,\ldots,1)\)，则 \(P\equiv I_m\pmod{q_i}\) 对所有 \(i\) 成立，但 \(\det P=M+1>1\)，并非整数环上的换基矩阵。
\end{solution}

\begin{exercise}\label{b2:ex:A-near-singular-solutions}
对 \(\varepsilon\ne0\)，精确求解
\[
\begin{pmatrix}1&1\\1&1+\varepsilon\end{pmatrix}
\binom{x}{y}=\binom{2}{2+\delta}.
\]
比较 \(\delta=0\) 与 \(\delta=\varepsilon\) 的解，说明当 \(\varepsilon\) 很小时，右端的小变化为什么仍可能改变解一个固定的量。若 \(\varepsilon\) 的输入仅为包含零的小区间，能否直接套用所得除法公式？
\end{exercise}
\begin{solution}
两行相减得到 \(\varepsilon y=\delta\)，因而
\[
y=\delta/\varepsilon,\qquad x=2-\delta/\varepsilon.
\]
\(\delta=0\) 时解为 \((2,0)\)，\(\delta=\varepsilon\) 时为 \((1,1)\)。右端只相差 \((0,\varepsilon)\)，解却相差 \((-1,1)\)。这是精确等式给出的比较，无须归因于某次浮点舍入。

如果输入区间包含零，就尚未知道 \(\varepsilon\ne0\)，不能统一除以它。在 \(\varepsilon=0\) 时，\(\delta=0\) 给出无穷多解，而 \(\delta\ne0\) 给出无解，与非零情形的唯一解不同。
\end{solution}

\begin{exercise}\label{b2:ex:A-polynomial-system}
对\eqref{b2:eq:A-polynomial-certificate}中的 \(B\in M_2\bigl(\mathbb F_5[t]\bigr)\)，判断方程
\[
Bx=\binom{0}{t^2},\qquad
Bx=\binom{1}{t-t^3+t^4}
\]
是否有多项式解；有解时求出全部解。第一式在分式域上是否可解？
\end{exercise}
\begin{solution}
第一式左乘 \(P\) 后右端仍为 \((0,t^2)^{\mathsf T}\)，而 \(t^4\nmid t^2\)，所以没有多项式解。第二式左乘 \(P\) 后右端为 \((1,t^4)^{\mathsf T}\)，对角方程的解是 \(y=(1,1)^{\mathsf T}\)。所以
\[
x=Qy=\binom{1-t}{1+t^2-t}.
\]
两个非零对角元在整环中都可消去，故齐次核为零，这也是唯一解。

第一式在 \(\mathbb F_5(t)\) 上可解，先得 \(y=(0,t^{-2})^{\mathsf T}\)，再乘 \(Q\) 得
\[
x=\binom{-t^{-1}}{1+t^{-2}}.
\]
它含有分母，不能当作原环上的多项式解。
\end{solution}

\begin{exercise}\label{b2:ex:A-rational-powers}
对\secref{b2:sec:A04}中的有理矩阵 \(T\)，求 \(T^{-1}\) 的各条目，并把 \(T^5\) 写成 \(I,T,T^2\) 的线性组合。用原矩阵验证 \(T^{-1}e_3=Tv\)。
\end{exercise}
\begin{solution}
直接乘法得到
\[
T^2=\begin{pmatrix}0&3&1\\0&2&2\\1&-\tfrac12&2\end{pmatrix}.
\]
由\eqref{b2:eq:A-rational-inverse}，
\[
T^{-1}=\begin{pmatrix}-1&\tfrac32&\tfrac12\\0&0&1\\\tfrac12&-\tfrac14&0\end{pmatrix}.
\]
又 \(T^3=2T+2I\)，所以 \(T^4=2T^2+2T\)，进而
\[
T^5=2T^3+2T^2=2T^2+4T+4I.
\]
逆矩阵的第三列为 \((1/2,1,0)^{\mathsf T}=Tv\)。原矩阵乘此列恰为 \(e_3\)，与 \(T^2v=e_3\) 一致。
\end{solution}

\begin{exercise}\label{b2:ex:A-equivalence-not-similarity}
令 \(A=\begin{pmatrix}1&1\\0&1\end{pmatrix}\)、\(B=I_2\)。证明它们在 \(\mathbb Z\) 上幺模等价，并且具有相同特征多项式，但在 \(\mathbb Q\) 上不相似。说明只用 Smith 因子判断有理矩阵的相似类会遗漏什么。
\end{exercise}
\begin{solution}
取 \(Q=\begin{pmatrix}1&-1\\0&1\end{pmatrix}\)，则 \(\det Q=1\)、\(AQ=I_2\)，两矩阵整数 Smith 因子都是 \(1,1\)。两者特征多项式也都是 \((t-1)^2\)。但 \(A-I_2\ne0\) 而 \((A-I_2)^2=0\)，故 \(A\) 的最小多项式为 \((t-1)^2\)；\(B\) 的最小多项式为 \(t-1\)。相似变换满足 \(f(S^{-1}AS)=S^{-1}f(A)S\)，保持零化多项式，故二者不相似。

矩阵自身的 Smith 因子描述其余核，不保存算子的反复作用。判断 \(A\) 在 \(\mathbb Q\) 上的相似类，应研究 \(\mathbb Q[t]\) 上的 \(tI-A\) 及相应不变因子，而不是只研究 \(A\) 在 \(\mathbb Q\) 或 \(\mathbb Z\) 上的行列等价。
\end{solution}

\begin{exercise}\label{b2:ex:A-specialization-loses-multiplicity}
令 \(R=\mathbb F_5[t]\)、\(M_r=R/(t^r)\)（\(r\ge1\)），并令 \(S_\ell=R/(t^\ell)\)（\(\ell\ge1\)）。求 \(M_r\otimes_R S_\ell\) 的 \(\mathbb F_5\) 维数。比较 \(M_4\) 与 \(M_8\) 在 \(\ell=1\) 和 \(\ell=5\) 时的结果，说明保留截断参数而不只取 \(t=0\) 后，怎样恢复被逐点秩遗漏的指数 \(r\)。
\end{exercise}
\begin{solution}
由\cref{b2:prop:tensor-quotient}，先对 \(M_r\) 商去 \(t^\ell M_r\)，得到
\[
M_r\otimes_R S_\ell
\cong R/(t^r,t^\ell)
\cong R/(t^{\min(r,\ell)}).
\]
最后一个商的基为 \(1,t,\ldots,t^{\min(r,\ell)-1}\) 的剩余类，维数为 \(\min(r,\ell)\)。\(\ell=1\) 正是取 \(t=0\) 的纤维，两个模都只给出一维；\(\ell=5\) 时，\(M_4\) 与 \(M_8\) 则分别给出四维和五维，已经可以区分。

随着 \(\ell\) 增大，这个维数序列依次为 \(1,2,\ldots,r,r,r,\ldots\)。因此其最终稳定值恰为 \(r\)。只取一个点会抹掉 \(t\) 的正次幂；在 \(S_\ell\) 中保留 \(t,t^2,\ldots,t^{\ell-1}\)，便可以逐次看到更高的重数信息。
\end{solution}
